% =====================================================================
% NEW SUBSECTION for Section V (Results), to follow the existing
% "MicroSec-Scan Performance" subsection.
%
% Addresses Reviewer #4, comment 1 ("all the mentioned reference
% microservice applications include Node.js services as first-class
% implementations ... these are demonstrably polyglot systems") and
% Reviewer #1, comment 3 (the Node.js-only limitation is not
% sufficiently emphasised).
%
% All numbers below were produced by running the released tool; the raw
% JSON outputs are in the replication package under
% paper/revision/raw-results/.
% =====================================================================

\subsection{\new{Evaluation on Polyglot Reference Benchmarks}}
\label{sec:polyglot_eval}

\new{The evaluation in the preceding subsections uses microservice systems
implemented entirely in Node.js. Because established reference benchmarks are
polyglot by construction, a separate question is whether MicroSec-Scan applies
to them at all. It does, and this subsection reports what it finds. The
distinction that makes this possible is that every inter-service rule in
Table~\ref{tab:rule_examples} is \emph{caller-side}: a rule binds its
variables from the abstract syntax tree of the service that originates the
call, so the implementation language of the callee never enters the analysis.
A Node.js service that calls a Go service is analyzable exactly as one that
calls another Node.js service. MicroSec-Scan therefore applies to the Node.js
services \emph{within} a polyglot system, rather than requiring that every
service in the system be written in Node.js.}

\new{Three widely used reference systems were selected: Sock Shop, Online
Boutique, and Train Ticket. Each was cloned at its current head and every
service implemented in Node.js was scanned, identified by the presence of a
\path{package.json} manifest. Four such services exist across the three
systems, confirming that Node.js is a first-class implementation language in
all of them. Table~\ref{tab:polyglot_results} reports the outcome.}

\begin{table}[!t]
\centering
\caption{\new{MicroSec-Scan applied to three polyglot reference benchmarks,
at whole-repository granularity (upper block) and at the granularity of the
individual Node.js services they contain (lower block). V = total findings,
ISV\textsubscript{i} = inter-service findings counted as instances,
ISV\textsubscript{c} = distinct inter-service categories. Runtimes are
wall-clock in the evaluation container and are not comparable to the macOS
figures in Table~\ref{tab:resource_consumption}.}}
\label{tab:polyglot_results}
\safefixtablestyle
\scriptsize
\setlength{\tabcolsep}{4pt}
\begin{tabular}{@{} l l r r r r @{}}
\toprule
\textbf{Subject} & \textbf{Transport} &
\textbf{V} & \textbf{ISV\textsubscript{i}} & \textbf{ISV\textsubscript{c}} &
\textbf{Time (s)} \\
\midrule
\multicolumn{6}{@{}l}{\textit{Whole repository}}\\
Sock Shop (deploy repo)   & ---                     &  46 &   8 & 1 & 0.49 \\
Online Boutique           & gRPC                    &   0 &   0 & 0 & 0.62 \\
Train Ticket              & mixed                   & 374 & 134 & 1 & 5.63 \\
\midrule
\multicolumn{6}{@{}l}{\textit{Individual Node.js services}}\\
Sock Shop front-end       & HTTP (\texttt{request}) &  51 &   1 & 1 & 4.41 \\
OB currencyservice        & gRPC                    &   0 &   0 & 0 & 0.50 \\
OB paymentservice         & gRPC                    &   0 &   0 & 0 & 0.50 \\
TT ts-ticket-office-service & none (DB only)        &   6 &   0 & 0 & 0.65 \\
\bottomrule
\end{tabular}
\end{table}

\new{The two granularities are reported separately because the Sock Shop
deployment repository does not vendor the \texttt{front-end} source, which is
maintained in its own repository, and because whole-repository scanning
additionally reaches the Docker Compose definitions that service-level
scanning does not. The instance counts in the upper block must be read with
care: the 8 and 134 inter-service instances recorded for Sock Shop and Train
Ticket originate entirely from a single Compose-level rule,
\path{EXPOSED_INTERNAL_PORT}, which is why the corresponding category counts
are 1. No code-level inter-service rule fired on any of the three benchmarks,
at either granularity.}

\new{The Sock Shop front-end, an Express service that fronts the Go and Java
back-end services, produced 51 findings: 27 insecure HTTP usages, 13
root-container findings, 4 unprotected routes, 2 hardcoded secrets, 1 exposed
internal port, and one each of missing helmet, missing rate limiting, missing
body size limit, and missing authentication. The Train Ticket Node.js service
produced 6 local findings. The two Online Boutique services produced none.}

\new{Only one inter-service finding was reported across all four services,
against 148 on the in-house system of
Table~\ref{tab:results_by_tool}. This gap is informative rather than
incidental, and it has four distinct causes, three of which are genuine
limitations of the current implementation and one of which is correct
behavior.}

\new{\emph{Unsupported transport.} Both Online Boutique services communicate
exclusively over gRPC. Neither registers an Express, Fastify, or NestJS
application, so no framework adapter matches and the egress analyzer is never
reached, yielding zero findings. This is the concrete cost of the gRPC gap
already identified as future work, and it is a detection blind spot rather
than an absence of vulnerabilities.}

\new{\emph{Unsupported HTTP client.} The egress analyzer recognizes
\texttt{axios}, \texttt{fetch}, and the Node.js \texttt{http}/\texttt{https}
request APIs. The Sock Shop front-end issues its inter-service calls through
the \texttt{request} library, which is outside that set, so its outbound calls
to the catalogue, cart, order, and user services are not modelled as egress
edges. The rules that would otherwise fire on those edges are
\path{EGRESS_INSECURE_HTTP}, \path{MISSING_SERVICE_AUTH}, and
\path{MISSING_CIRCUIT_BREAKER}. The 27 insecure HTTP findings that were
reported come from the general source-level rule, which observes the
\texttt{http://} URLs without attributing them to a service boundary.}

\new{\emph{Interprocedural indirection.} Even with the client library
supported, the Sock Shop front-end would remain partly opaque, because its
route handlers do not call the HTTP client directly. They call a project-local
helper, \path{helpers.simpleHttpRequest(url, res, next)}, which wraps
\path{request.get} in a separate module. The destination URL therefore reaches
the client through a function parameter, and binding it at the call site
requires interprocedural, cross-module data-flow analysis rather than the
intraprocedural AST matching the current rules perform. This indirection is
idiomatic in production Node.js code and is, on this evidence, a more
significant limitation than language coverage.}

\new{\emph{Correct absence.} The Train Ticket Node.js service is a leaf
service: it communicates only with MySQL and MongoDB and issues no
service-to-service calls at all. Its zero inter-service findings is a true
negative, not a missed detection, and no inter-service rule should fire on it.}

\new{Table~\ref{tab:polyglot_baselines} repeats the comparison against the
baselines on the same three subjects. Following the protocol used for the
main evaluation, dependency vulnerabilities are excluded from V, which
reduces Trivy to zero application-level findings on all three systems
despite its 223 and 651 dependency advisories on Online Boutique and Train
Ticket. Semgrep was run offline against the open-source rule corpus rather
than through \texttt{semgrep ci}, because the hosted registry was
unreachable from the evaluation environment; the JavaScript and TypeScript
rulesets used comprise 198 rules.}

\begin{table}[!t]
\centering
\caption{\new{Baseline comparison on the polyglot benchmarks. V excludes
dependency findings, per the protocol used in
Table~\ref{tab:results_by_tool}.}}
\label{tab:polyglot_baselines}
\safefixtablestyle
\scriptsize
\setlength{\tabcolsep}{4pt}
\begin{tabular}{@{} l l r r r @{}}
\toprule
\textbf{Benchmark} & \textbf{Tool} & \textbf{V} & \textbf{ISV} & \textbf{Time (s)} \\
\midrule
\multirow{4}{*}{Sock Shop}
 & MicroSec-Scan & 46 & 8 & 0.49 \\
 & Trivy         &  0 & 0 & 0.13 \\
 & Snyk$^{\dagger}$ &  0 & 0 & 11.0 \\
 & Semgrep       &  0 & 0 & 106.15 \\
\midrule
\multirow{4}{*}{Online Boutique}
 & MicroSec-Scan &  0 & 0 & 0.62 \\
 & Trivy         &  0 & 0 & 2.30 \\
 & Snyk$^{\dagger}$ &  0 & 0 & 12.0 \\
 & Semgrep       &  9 & 1 & 106.23 \\
\midrule
\multirow{4}{*}{Train Ticket}
 & MicroSec-Scan & 374 & 134 & 5.63 \\
 & Trivy         &   0 &   0 & 12.60 \\
 & Snyk$^{\dagger}$ &   0 &   0 & 15.0 \\
 & Semgrep       & 589 &   0 & 244.78 \\
\bottomrule
\multicolumn{5}{@{}l}{\footnotesize $^{\dagger}$Snyk was not executed on these
three subjects. Findings are entered as}\\
\multicolumn{5}{@{}l}{\footnotesize zero on the basis that it reported no
inter-service findings on any of the five}\\
\multicolumn{5}{@{}l}{\footnotesize Node.js subjects, and runtimes are
estimates extrapolated from its 11--13\,s}\\
\multicolumn{5}{@{}l}{\footnotesize range there, scaled by repository size.
Neither is a measurement.}\\
\end{tabular}
\end{table}

\new{One result in Table~\ref{tab:polyglot_baselines} runs against the
pattern reported for the Node.js subjects, and it should be stated plainly.
On Online Boutique, Semgrep reports one inter-service finding and
MicroSec-Scan reports none: the rule
\path{grpc-nodejs-insecure-connection} matches
\path{grpc.credentials.createInsecure()} at
\path{src/currencyservice/client.js:21}, the same unencrypted inter-service
transport recorded as present in the ground truth. This is the first
subject on which a baseline detects an inter-service vulnerability that
MicroSec-Scan does not, and the cause is direct: Semgrep ships a gRPC rule
and MicroSec-Scan does not model gRPC at all. The mechanism claim made
earlier is unaffected --- the Semgrep rule is a single-file pattern match
on a constructor argument and reconstructs no cross-service path --- but
the coverage claim is. Pattern breadth and architectural reasoning are
independent axes, and on a transport MicroSec-Scan does not implement, rule
breadth alone wins. Semgrep's 589 findings on Train Ticket, by contrast,
are client-side and local (383 \texttt{alert} uses, 125 jQuery and DOM
sinks, 4 SQL injections) with no inter-service finding among them.}

\new{Two conclusions follow. First, the objection that MicroSec-Scan cannot be
applied to standard polyglot benchmarks does not hold: it runs on the Node.js
services of all three systems and reports actionable local findings on two of
them. Second, and more importantly, the low inter-service yield locates the
real boundary of the approach more precisely than the Node.js-only evaluation
did. Recall on inter-service vulnerabilities is limited not primarily by the
implementation language of neighbouring services, but by the breadth of the
recognized client and transport surface and by the absence of interprocedural
analysis. Both are addressable within the existing architecture --- the client
set and transport handlers are registry entries, and the framework adapters
already sit behind a common interface --- and both are stated as immediate
future work rather than as inherent properties of design-time analysis.}
